% Brief (tikz-diagram-craft).
% Reader question: given four overlapping requests, what is the one
%   commit order the single writer actually produces?
% Claim: a concrete three-agent race -- A claims, B loses and is told who
%   holds it, A releases, C claims -- linearizes to one total order on the
%   writer's rail, even though the requests overlap in caller time.
% Evidence: whitepaper/single-writer-kernel.tex thm:consistency (the
%   conditional consistency model: committed transactions admit a serial
%   order, claim/lock operations are linearizable with commit as the
%   linearization point).
% Must distinguish: caller-time overlap (the bands) from the one commit
%   order (the rail); that a losing claim still gets a definite outcome
%   (told who holds it), not a stall.
% Grammar: swimlane bands over a shared commit rail, the worked instance's
%   own linearization written under it (S13's K&R register: a concrete
%   record, not a schematic). Rejected: a formal transition-system diagram,
%   which would state the theorem again rather than show one witness to it.
% Five-second test: "three overlapping bars, one rail, four commit points in
%   order, and the one linearization written out".
\begin{figure}[H]
\centering
\pdfigurehue{pdcobalt}%
\begin{tikzpicture}[pd figure,x=.8cm,y=.85cm]
  \foreach \y/\a in {5.4/A,4.4/B,3.4/C} {
    \node[pd title,anchor=east] at (-.25,\y) {agent \a};
    \draw[pd hairline] (0,\y)--(10.4,\y);
  }
  \fill[pd neutral fill] (.6,5.2) rectangle (4.3,5.6);
  \node[pd label,anchor=south,text width=1.2cm,align=center] at (2.45,5.7) {claim};
  \fill[pd neutral fill] (1.9,4.2) rectangle (6.0,4.6);
  \node[pd label,anchor=north,text width=1.9cm,align=center] at (2.9,4.10) {loses to A};
  \fill[pd neutral fill] (6.4,5.2) rectangle (8.0,5.6);
  \node[pd label,anchor=south,text width=1.1cm,align=center] at (7.2,5.7) {release};
  \fill[pd neutral fill] (8.4,3.2) rectangle (10.3,3.6);
  \node[pd label,anchor=south,text width=.9cm,align=center] at (9.35,3.7) {wins};

  \foreach \x/\lbl in {4.3/{A claim},6.0/{B loses},8.0/{A release},10.3/{C claim}} {
    \draw[pd guide] (\x,1.55)--(\x,5.75);
    \node[pd focus datum] at (\x,1.9) {};
    \node[pd label,anchor=north,align=center] at (\x,1.65) {\lbl};
  }
  \draw[pd focus arrow] (1.1,1.9)--(10.5,1.9);
  \node[pd title,anchor=east,text width=1.6cm,align=right] at (.7,1.9) {commit order};

  \node[pd note,anchor=north,text width=10.4cm,align=center] at (5.2,.75)
    {$A.claim \prec B.claim\text{ (loses)} \prec A.release \prec C.claim$};
\end{tikzpicture}
\caption{Overlapping requests receive one commit order: A claims, B loses,
A releases, C claims. Theorem~\ref{thm:consistency} states the five
assumptions this ordering requires. [internal]}
\label{fig:swk-consistency-model}
\end{figure}
